\newpage
\section{The Multi-Prover Verification Framework}

\subsection{Overview of \texttt{self\_proving.py}}

The \texttt{self\_proving.py} module extends Layer~1 with formal proof capabilities. It provides the \texttt{SelfProvingAtomicity} class, which is attached to an \texttt{UltimateObservable} instance and exposes a single public method:

\begin{lstlisting}[language=Python, caption={Public interface of \texttt{SelfProvingAtomicity}.}]
def prove_atomicity(self, framework: str,
                    provers: Optional[List[TheoremProverType]] = None,
                    timeout_ms: int = 5000) -> Optional[Proof]
\end{lstlisting}

This method attempts to prove that the observable is atomic according to the specified framework (\texttt{"boolean"}, \texttt{"measure"}, \texttt{"category"}, or \texttt{"info"}) using a configurable list of theorem provers. It returns a \texttt{Proof} object containing the generated proof artifacts and verification status.

Internally, the module contains:
\begin{itemize}
    \item \texttt{AtomicityTheoremGenerator}: Translates the abstract atomicity claim into a formal statement suitable for each prover.
    \item \textbf{Prover-specific backends:} Functions that invoke SymPy, Z3, Lean~4, and Coq, capture their output, and update the \texttt{Proof} object.
    \item \textbf{Certificate generation:} The \texttt{Proof.to\_certificate()} method exports a JSON document containing the proof statement, verification status, and (for SymPy) the simplified logical formula.
\end{itemize}

% figures/architecture_overview.tex
\begin{figure}[htbp]
\centering
\begin{tikzpicture}[
    node distance = 1.2cm and 1.8cm,
    box/.style = {draw, rounded corners, minimum width=3.0cm, minimum height=0.9cm, align=center, fill=blue!5, font=\small},
    prover/.style = {draw, rounded corners, minimum width=2.2cm, minimum height=0.8cm, align=center, fill=green!5, font=\small},
    arrow/.style = {-{Stealth[length=3mm]}, thick},
    every label/.style = {font=\small\ttfamily}
]

% Top layer: UltimateObservable
\node[box, fill=gray!15] (obs) {\texttt{UltimateObservable}\\\small (value, meta\_spec)};

% SelfProvingAtomicity
\node[box, below=of obs] (prover) {\texttt{SelfProvingAtomicity}\\\small (cache, lock)};

% Theorem Generator
\node[box, below left=of prover, xshift=-0.8cm] (gen) {\texttt{AtomicityTheoremGenerator}\\\small (statement, formula)};

% Provers matrix
\matrix[below=of prover, column sep=0.4cm, row sep=0.3cm, xshift=1.0cm] (provers) {
    \node[prover] (sympy) {SymPy\\(logical simplification)}; &
    \node[prover] (z3) {Z3\\(SMT solver)}; \\
    \node[prover, yshift=-0.5cm] (lean) {Lean 4\\(interactive)}; &
    \node[prover, yshift=-0.5cm] (coq) {Coq\\(interactive)}; \\
};

% Proof object
\node[box, fill=yellow!10, below=of gen, yshift=-1.5cm] (proof) {\texttt{Proof}\\\small (statement, verified\_by, cert)};

% Arrows
\draw[arrow] (obs.south) -- (prover.north);
\draw[arrow] (prover.south) -- ++(0,-0.3) -| (gen.north);
\draw[arrow] (prover.south) -- ++(0,-0.3) -| (provers.north);

\draw[arrow] (gen.south) -- ++(0,-0.3) -| (proof.north);
\draw[arrow] (sympy.south) -- ++(0,-0.3) -| (proof.north);
\draw[arrow] (z3.south) -- ++(0,-0.3) -| (proof.north);
\draw[arrow] (lean.south) -- ++(0,-0.3) -| (proof.north);
\draw[arrow] (coq.south) -- ++(0,-0.3) -| (proof.north);

% Feedback loop (cache)
\draw[arrow, dashed, bend right=25] (proof.east) to node[above, sloped, font=\footnotesize] {cache} (prover.south west);

\end{tikzpicture}
\caption{Architecture of the multi-prover verification module. The \texttt{SelfProvingAtomicity} class manages the proof process, using the \texttt{AtomicityTheoremGenerator} to produce formal statements and dispatching them to the configured provers. Results are aggregated in a \texttt{Proof} object and cached for future queries.}
\label{fig:architecture}
\end{figure}
\subsection{Translation Semantics and Prover Encodings}

A central challenge in multi-prover verification is ensuring that the property being proved is semantically preserved across fundamentally different logical systems. In this section, we formally specify the atomicity properties for each framework and detail their concrete encoding in SymPy, Z3, Lean~4, and Coq. We conclude with a critical discussion of the trust assumptions inherent in our translation layer.

\subsubsection{Formal Specification of Multi-Axial Atomicity}

We first restate the four atomicity criteria from the \Apeiron{} Framework in precise mathematical language, which serves as the \textit{specification} that all prover backends must implement.

\begin{enumerate}[label=(\roman*), leftmargin=*]
    \item \textbf{Boolean Atomicity:} For an observable $x$, let $B$ be a Boolean algebra with partial order $\leq$. $x$ is a Boolean atom iff $x \neq \bot$ and there exists no $y \in B$ such that $\bot < y < x$.
    \item \textbf{Measure-Theoretic Atomicity:} For a set $A$ and a measure $\mu$, $A$ is a measure atom iff $\mu(A) > 0$ and for every measurable $B \subseteq A$, either $\mu(B) = 0$ or $\mu(B) = \mu(A)$.
    \item \textbf{Categorical Atomicity:} For an object $X$ in a category $\mathcal{C}$, $X$ is a zero object (and thus atomic) iff it is both initial (unique morphism to every object) and terminal (unique morphism from every object).
    \item \textbf{Information-Theoretic Atomicity:} For a string $s$, let $K(s)$ denote its Kolmogorov complexity. $s$ is information-theoretically atomic iff $K(s) \approx |s|$; i.e., it is algorithmically incompressible.
\end{enumerate}

In practice, these abstract definitions are instantiated on the concrete value stored in an \texttt{UltimateObservable} (e.g., an integer, a set, a string). The \texttt{AtomicityTheoremGenerator} translates these value-specific claims into prover-native representations.

\subsubsection{Encoding Across Provers}

Table~\ref{tab:prover_encodings} summarizes how each atomicity framework is encoded for the four supported provers. The encodings reflect the unique strengths and idiomatic constructs of each system.

\begin{table}[htbp]
\centering
\caption{Concrete encodings of atomicity properties for each prover.}
\label{tab:prover_encodings}
\small
\begin{tabular}{p{2.2cm} p{2.8cm} p{2.8cm} p{2.5cm} p{2.5cm}}
\toprule
\textbf{Framework} & \textbf{SymPy} & \textbf{Z3} & \textbf{Lean 4} & \textbf{Coq} \\
\midrule
Boolean (int) & \texttt{And(atomic, Not(decomposable))} & \texttt{Not(Exists(x, And(x>1, x<val, val\%x==0)))} & \texttt{isAtom val} (via divisor search) & \texttt{is\_prime} (for primes) \\
Measure (set) & \texttt{len(set)==1} & \texttt{Length(set) == 1} & \texttt{cardinality} equality & \texttt{length} property \\
Category & \texttt{And(initial, terminal)} & \texttt{initial \&\& terminal} & Zero object construction & Trivial (placeholder) \\
Information & \texttt{compress\_ratio >= 0.85} & \texttt{compressed\_len >= 0.85 * original\_len} & Trivial (placeholder) & Trivial (placeholder) \\
\bottomrule
\end{tabular}
\end{table}

\paragraph{SymPy.} The generator constructs Boolean expressions using SymPy's symbolic primitives (\texttt{Symbol}, \texttt{And}, \texttt{Not}, \texttt{Implies}). For information-theoretic atomicity, it computes the zlib compression ratio and creates an expression asserting the ratio exceeds a threshold. Verification succeeds if \texttt{simplify\_logic} does not reduce the formula to \texttt{False}.

\paragraph{Z3.} Z3's SMT-LIB language supports quantifiers and arithmetic, enabling direct encoding of the existential conditions. For Boolean atomicity, the generator emits a formula asserting the non-existence of a proper divisor. For measure atomicity, it asserts that the set has exactly one element. The solver checks for satisfiability of the negated claim; \texttt{unsat} indicates a successful proof.

\paragraph{Lean 4 and Coq.} For these interactive proof assistants, the generator produces a theorem statement and a proof script. Currently, non-trivial proofs are only generated for Boolean atomicity of integers; for other frameworks, a trivial proof (\texttt{trivial} in Lean, \texttt{Proof. trivial. Qed.} in Coq) is used as a placeholder. Fully fleshed-out formalizations of measure and category theory in Lean's mathlib or Coq's standard library are planned for future work.

\subsubsection{Limitations of the Translation Layer}

The current implementation employs \textit{string templating} to generate prover-specific code. While pragmatic, this approach introduces several trust and correctness concerns:

\begin{itemize}
    \item \textbf{No formal guarantee of semantic preservation:} There is no proof that the generated Lean theorem \texttt{isAtom 1} corresponds exactly to the Python function \texttt{boolean\_atomicity(1)}. The translation relies on the programmer's manual alignment of definitions.
    \item \textbf{Syntactic fragility:} Observable IDs containing special characters (e.g., backticks, quotes) can break the generated Lean or Coq syntax. Basic sanitization is performed, but a robust parser is lacking.
    \item \textbf{Incomplete coverage:} Non-Boolean frameworks currently use placeholder proofs in Lean and Coq, meaning that cross-verification for those frameworks is only as strong as SymPy and Z3.
    \item \textbf{Lack of error recovery:} If a prover fails due to a malformed script, the system simply marks that prover as unsuccessful; it cannot diagnose or correct the translation error.
\end{itemize}

These limitations are intrinsic to the current rapid-prototyping stage of the \Apeiron{} verification module. For production-grade formal assurance, the translation layer should be rebuilt around an \textbf{Abstract Syntax Tree (AST)} in Python, with formally verified serializers for each target language. The multi-prover consensus mechanism described in Section~3 provides a pragmatic, albeit heuristic, safeguard against individual translation errors.

In the next section, we detail the theorem generation process, which builds upon the encodings described here.

\subsection{Theorem Generation}

For each framework, the \texttt{AtomicityTheoremGenerator} produces a human-readable theorem statement and, where applicable, a formal representation in the prover's native language.

\subsubsection{SymPy (Logical Simplification)}
For Boolean and information-theoretic atomicity, the generator constructs a SymPy Boolean expression. For example, for Boolean atomicity of an integer $n$:
\begin{lstlisting}[language=Python]
atomic = Symbol('atomic')
decomposable = Symbol('decomposable')
base = And(Implies(atomic, Not(decomposable)),
           Implies(Not(decomposable), atomic))
# Additional axioms specific to the integer value are conjoined.
\end{lstlisting}
SymPy's \texttt{simplify\_logic} is then used to reduce the formula. If the simplified result is not \texttt{False}, the proof is considered successful (i.e., the negation of atomicity is not forced).

\subsubsection{Z3 (SMT Solving)}
For Z3, the generator creates a Z3 Boolean expression. For Boolean atomicity, it directly encodes the definition: an integer $n$ is a Boolean atom if no proper divisor exists.
\begin{lstlisting}[language=Python]
x = Int('divisor')
solver = Solver()
solver.add(x > 1, x < n, n % x == 0)
is_atom = (solver.check() == unsat)
\end{lstlisting}
For measure-theoretic atomicity, it checks that the set has exactly one element. For information-theoretic atomicity, it checks the incompressibility ratio of the byte representation.

\subsubsection{Lean~4 and Coq (Interactive Proof Assistants)}
For Lean~4, the generator produces a Lean theorem statement and a proof script. For Boolean atomicity of the integer~1, the generated Lean code is:
\begin{lstlisting}[language=Python]
theorem atomicity_obs : isAtom 1 := by
  constructor
  - decide
  - intro b h1 h2; omega
\end{lstlisting}
For Coq, a trivial proof (\texttt{Proof. trivial. Qed.}) is generated for non-Boolean frameworks; a more substantive proof could be provided for the Boolean case.

\subsection{Prover Invocation and Verification}

Each prover backend executes the generated proof script in a subprocess or via API calls, with a configurable timeout. The results are captured and used to update the \texttt{Proof} object:
\begin{itemize}
    \item \textbf{SymPy:} The simplified formula is stored in \texttt{proof.proof\_sympy}. If the simplification does not yield \texttt{False}, \texttt{proof.add\_verification("sympy", elapsed\_ms)} is called.
    \item \textbf{Z3:} If \texttt{solver.check()} returns \texttt{unsat}, the proof is considered verified, and the SMT-LIB representation is stored in \texttt{proof.proof\_z3}.
    \item \textbf{Lean~4:} A temporary \texttt{.lean} file is written and executed with \texttt{lean --run}. A return code of $0$ indicates success, and the proof script is stored in \texttt{proof.proof\_lean}.
    \item \textbf{Coq:} A temporary \texttt{.v} file is compiled with \texttt{coqc}. A return code of $0$ indicates success, and the proof script is stored in \texttt{proof.proof\_coq}.
\end{itemize}

The \texttt{Proof} object maintains a list \texttt{verified\_by} containing the names of all provers that successfully discharged the proof. The overall \texttt{status} becomes \texttt{ProofStatus.VERIFIED} if at least one prover succeeded.

\subsection{Certificate Export}

The \texttt{Proof.to\_certificate()} method serializes the proof into a JSON structure:
\begin{lstlisting}[caption={Example certificate for Boolean atomicity of integer~1.}]
{
  "statement": "Theorem atomicity_one_boolean: ...",
  "status": "verified",
  "verified_by": ["sympy", "z3", "lean", "coq"],
  "verification_time_ms": 10.3,
  "created_at": 1776710478.8289564,
  "has_lean": true,
  "has_coq": true,
  "has_sympy": true,
  "has_z3": true,
  "proof_sympy": "atomic_one_boolean & ~decomposable_one_boolean",
  "metadata": {"framework": "boolean"}
}
\end{lstlisting}
This certificate can be stored, shared, and independently re-verified using the same (or compatible) provers. The inclusion of the SymPy proof string allows for lightweight re-checking without requiring the original Python environment.

\subsection{Caching and Thread Safety}

All proof computations are cached per framework and prover combination. The cache is stored in the \texttt{SelfProvingAtomicity} instance attached to the observable, and access is protected by a reentrant lock (\texttt{threading.RLock}). This ensures that repeated queries for the same proof are served instantly, which is essential for interactive applications and higher-layer algorithms that may request atomicity proofs multiple times.